\section{Broader Impact \& Limitations}
\label{appendix:limitations}

Currently, large language models (LLMs) face two typical challenges stemming from the Transformer architecture: limited context window capacity and the hallucination problem. Retrieval-augmented techniques offer a viable solution to these critical issues and have gained popularity in various application scenarios. The proposed \ours~ achieves high-precision retrieval in diverse complex scenarios, thereby driving comprehensive improvements in human-computer interaction. It enables AI systems to provide more accurate and stable services to society, facilitating the broader integration of AI technologies into daily life.

However, \ours~ still exhibits the following limitations: (1) For extremely large documents, constructing a knowledge graph requires substantial time and computational resources; however, it is still much faster than GraphRAG. (2) While the initial retrieval process with LLM-guided knowledge graph traversal leads to significant accuracy improvements, substantial computational resource consumption remains. Our efficiency enhancements focus on subsequent searches; however, further optimization is needed to address the high computational overhead resulting from multiple LLM calls during initialization.